\project{11}{An AT engine that never blocks}{NUCLEO-H7A3ZI-Q with the Waveshare SIM7020E}{Asynchronous command queue, unsolicited results, timeouts, backoff}

\begin{keyfacts}
\item[Board] NUCLEO-H7A3ZI-Q, with the Waveshare SIM7020E narrow-band modem on its own supply
\item[Peripherals] A USART with circular DMA receive and idle-line detection, a general timer for the timeout wheel, GPIO for the module's power control line
\item[Toolchain] arm-none-eabi-gcc with CMake on the target; a Python fault injector on the host
\item[Operating system] Bare metal, which is the whole point of the chapter
\item[Difficulty] 4 of 5
\item[Effort] 4 evenings of about four hours
\item[Deliverable] A command engine whose main loop never waits, which survives one hundred injected faults with every one ending in a defined state, and which stops in a named state when no SIM card is fitted
\end{keyfacts}

\subsection*{Why this project}

The usual driver for a modem of this kind is a chain of functions that each
write a command and then sit in a loop waiting for a reply. It works on a bench
with a good signal and it fails everywhere else, because everything the module
does that is not an immediate reply looks to that code like a stopped system. A
reply that arrives after the code gave up is discarded, and then answers the
next command. A message the module volunteers on its own, which is ordinary
behaviour and not an error, is read as the reply to whatever command happens to
be outstanding. A network registration that takes tens of seconds blocks the
sensor loop for tens of seconds. None of this is exotic; it is the normal life
of a cellular module, and a blocking driver has no vocabulary for it.

This chapter builds the alternative, which is small enough to hold in your head:
a queue of command records, one command in flight at a time, a per-command
timeout taken from the module's own documentation rather than guessed, a table
of unsolicited result codes dispatched independently of whatever command is
running, and an exponential backoff with a cap for the cases where retrying is
the right answer. The main loop calls one function that returns immediately, and
every path through the engine ends somewhere named.

The reason this is a chapter rather than a dependency is a gap in the prior art
that is worth stating precisely. Majerle's cellular library is MIT licensed, is
well engineered, explicitly supports the SIM7070G that is also on this bench, and
even bundles a messaging client. It requires an operating system and has no
bare-metal path. That is exactly the gap this chapter fills, and it is also the
reason chapter 20 exists: once there is a kernel, that library becomes the right
answer and this engine becomes the thing you understand it with. The same
library does not list the SIM7600E-H or the SIM7020E at all, and the silicon
vendor's own cellular middleware covers none of the three modems on this bench,
targeting a different set of module vendors entirely. There is no maintained
permissive component that does this job for this module on bare metal.

\begin{note}{The narrow-band module has no point-to-point protocol}
On the SIM7600E-H there is a choice: run a point-to-point link and put a network
stack on the microcontroller, or drive the module by commands. The SIM7020E
removes the choice. It offers no point-to-point protocol, so every byte of
payload is handed over inside a command and every result comes back as text.
That is not a restriction to design around; it decides the architecture, and it
is why this is necessarily a command-set chapter.
\end{note}

\subsection*{Prior art and what to reuse}

\begin{table}[H]\centering\small
\begin{tabular}{p{34mm}p{46mm}p{38mm}p{20mm}}
\toprule
Source & What it gives & What it does not & Licence \\
\midrule
Majerle's cellular library & A complete, maintained engine with a command queue, unsolicited result handling and a bundled messaging client, and it explicitly supports the SIM7070G that is also on this bench & It requires an operating system and has no bare-metal path, and it does not list the SIM7600E-H or the SIM7020E & MIT \\
The silicon vendor's cellular middleware & Nothing applicable here & It covers none of the three modems on this bench, targeting a different set of module vendors & Vendor licence \\
Majerle's host-side parser for the radio coprocessor & The clearest readable model of how a command queue, a line classifier and an unsolicited dispatcher fit together & It speaks a different command set, so the structure transfers and the commands do not & MIT \\
A peer-reviewed empirical study of this exact module & Realistic expectations for signal level, registration time and failure rate, which is what the retry and backoff policy has to be sized against & Its host is a single-board computer running a general-purpose operating system, not a bare-metal microcontroller & Paper \\
Chapters 2 and 4 of this volume & The byte path: a ring buffer and a circular DMA receive with idle-line detection, which is what makes a non-blocking engine possible at all & No framing of replies and no notion of a command & This volume \\
\bottomrule
\end{tabular}
\caption{Prior art for chapter 11. The one library that would do this job needs a kernel, which is the gap this chapter fills and the argument chapter 20 answers from the other side.}
\end{table}

What is left to write is the engine itself: the command record, the queue, the
timeout wheel, the line classifier, the unsolicited dispatcher, the backoff
policy and the fault injector that proves the whole thing. The byte path below
it is chapters 2 and 4, unchanged. The structural model is borrowed from the
host-side parser named above, which is the clearest published statement of how
these pieces fit together even though it speaks to different hardware.

One more thing is borrowed and it is the most important: the module's own AT
command manual states a maximum response time for each command. Those numbers
are the only defensible source for a per-command timeout. A timeout invented at
the keyboard is a number that will be wrong in the field in one direction or the
other, and the direction that fails quietly is the one that is too long.

\subsection*{Parts from the inventory}

\begin{table}[H]\centering\small
\begin{tabular}{p{40mm}p{66mm}p{40mm}}
\toprule
Part & Role & Interface \\
\midrule
NUCLEO-H7A3ZI-Q & Runs the engine & Micro USB to the host for flashing and the console \\
Waveshare SIM7020E & The narrow-band module under test. No point-to-point protocol, so command set only & Serial at 3.3 V logic, plus a power control line \\
Joy-it SBC-POW-BB or a USB supply & The module's own supply. It is never fed from the board & Barrel or USB, with the ground made common to the board \\
Renkforce USB/TTL cable & Talking to the module directly from the host, before any firmware exists, and again when the firmware disagrees with the module & Serial at 3.3 V. Disconnect the red 5 V lead \\
Waveshare SIM7070G & The comparison module. The library named above supports it, which makes it the control case & Serial \\
Host PC & Terminal, then the fault injector & USB \\
\bottomrule
\end{tabular}
\caption{Inventory items used in chapter 11. A SIM card is not in the inventory list, and the acceptance test is written so that its absence is a pass rather than an obstacle.}
\end{table}

\subsection*{System architecture}

\diagram{c11_arch}{The engine sits between an application that never waits and a byte path that never parses. Each layer has one job, and the two most common designs for this problem both come from collapsing two of these layers into one.}

Two collapses are worth naming because both are common and both are expensive.
Parsing inside the receive interrupt puts an unbounded amount of work at high
priority and makes every later timing measurement meaningless. Letting the
application call a send-and-wait function puts the modem's latency into the
sensor loop, which is the problem this chapter exists to remove. The layering in
the figure is not ceremony; each boundary is there because of a failure that
happens without it.

\subsection*{Peripheral configuration}

\begin{table}[H]\centering\small
\begin{tabular}{p{22mm}p{26mm}p{24mm}p{30mm}p{30mm}}
\toprule
Peripheral & Mode & Clock source & Pins and function & Interrupt and transfers \\
\midrule
USART to modem & Asynchronous, 115200 8N1 to start, module default confirmed first & Peripheral bus & A free USART on the Zio header. Which instance reaches Arduino D0 and D1 is on the confirm list & Circular DMA receive plus idle line; DMA transmit \\
USART3 console & Asynchronous, 115200 & Peripheral bus & Confirm in the board manual & Interrupt driven, chapter 3 \\
General timer & Up counter, 1 kHz tick & Peripheral bus & None & Update interrupt drives the timeout wheel \\
GPIO power control & Push-pull output & Peripheral bus & One free Zio pin & None \\
GPIO status input & Input with pull, optional & Peripheral bus & One free Zio pin & None. Polled \\
\bottomrule
\end{tabular}
\caption{Peripheral configuration for chapter 11. The baud rate is confirmed against the module before it is written into a header, because a module that has been left in a different rate by a previous experiment answers nothing and looks like a dead module.}
\end{table}

The timeout wheel does not need a dedicated timer. A millisecond tick from any
source works, and on a part with fifteen 16-bit timers and two 32-bit ones there
is no shortage. What it does need is to be independent of the command path, so
that a command which never completes still has something counting on its behalf.

\subsection*{Wiring}

\diagram{c11_wiring}{Two supplies, one ground. The module's peak demand is the reason for the first, and the fact that serial signalling is referenced to ground is the reason for the second. The cross between transmit and receive is the single most common wiring fault on this bench.}

\begin{note}{The module takes 2 A peaks and never takes them from the board}
A cellular module draws short, high current peaks when its transmitter turns on.
The figure on this bench is 2 A, which is more than the board can supply and
more than the measurement instrument of chapters 7 and 10 will pass. The module
gets its own supply, the grounds are made common so the serial signalling has a
reference, and the power instrument stays out of this path entirely. If the
module resets whenever it tries to register, the supply is the first suspect and
usually the only one.
\end{note}

A second safety point applies to the USB to serial cable used for talking to the
module directly. Disconnect its red 5 V lead. The module has its own supply and
a second one arriving down the serial cable is at best redundant.

\subsection*{Memory and timing budget}

\diagram{c11_mem}{Where the engine's bytes live. Everything is statically allocated, because an engine whose failure mode is exhaustion of a heap has replaced one unbounded behaviour with another.}

\begin{table}[H]\centering\small
\begin{tabular}{p{44mm}p{28mm}p{28mm}p{30mm}}
\toprule
Quantity & Budget & Measured & Margin \\
\midrule
Engine code in flash & 6 kB & not measured & not measured \\
Command queue, 8 slots & 512 B & not measured & not measured \\
Receive ring & 1 kB & not measured & not measured \\
Line assembly buffer & 256 B & not measured & not measured \\
Unsolicited result table & 256 B & not measured & not measured \\
Stack used by the engine & 512 B & not measured & not measured \\
Time in the engine per tick & below 20 \textmu{}s & not measured & not measured \\
Longest path with interrupts off & below 5 \textmu{}s & not measured & not measured \\
Commands faulted without a hang & 100 of 100 & not measured & not measured \\
\bottomrule
\end{tabular}
\caption{The budget for chapter 11. The two timing rows are what "never blocks" means as a number, and they are measured with the cycle counter rather than asserted. Nothing enters the measured column until the bench produces it.}
\end{table}

The receive ring is sized for the longest burst the module can produce between
two visits to the engine, not for the longest reply. Those are different
quantities and confusing them is how a ring that looks generous overflows. The
longest reply matters for the line assembly buffer instead, and a line longer
than that buffer is truncated deliberately and reported as a truncated line
rather than silently joined to the next one.

\subsection*{Firmware design (UML)}

\diagram{c11_uml}{The engine as a state machine. Every arrow out of a waiting state is either a reply, a timeout or an unsolicited message, and there is no arrow that leads nowhere. The two stopped states are outcomes and not faults.}

The state machine has one property that carries most of its value: the states
where the engine waits are the states where the application is free. In the
waiting states the engine holds a deadline and nothing else, and the function
the main loop calls returns immediately after checking whether the deadline has
passed and whether a complete line has arrived. Nothing in the engine ever spins.

Unsolicited messages are handled outside the state machine on purpose. A message
the module volunteers is not a reply and must not be treated as one, even when
one is outstanding. The line classifier decides which it is before the state
machine sees it, using the simple and reliable rule that a final result code
ends a command and anything matching an entry in the unsolicited table is
dispatched to that entry's handler no matter what state the engine is in.

\subsection*{Data flow (ASCII)}

\begin{asciiart}
  application           AT engine                     byte path             module
  +---------------+    +-------------------------+   +----------------+    +----------+
  | sensor loop   |    | queue: 8 command slots  |   | DMA tx         |    |          |
  |   |           |    |   +---------------+     |-->| USART tx       |--->|          |
  |   | at_submit |--->|   | cmd, timeout, |     |   +----------------+    | SIM7020E |
  |   |           |    |   | retries, done |     |   +----------------+    |          |
  |   | at_poll   |--->|   +---------------+     |   | DMA rx, circular|<---|          |
  |   |  returns  |    | one in flight at a time |<--| idle line        |   +----------+
  |   |  at once  |    |   deadline from the     |   | ring buffer      |       |
  |   v           |    |   module's own manual   |   +----------------+        | own
  | done callback |<---+-------------------------+          |                  | 2 A
  +---------------+    | line classifier              <-----+                  | supply
                       |   final result -> completes the command in flight
                       |   intermediate  -> appended to that command's result
                       |   unsolicited   -> dispatched from the table, any state
                       +-------------------------------------------------------+
\end{asciiart}

\subsection*{Repository layout}

\begin{plaincode}
nucleo-h7a3-at-engine/
  CMakeLists.txt
  src/at_engine.c              # queue, state machine, timeout wheel
  src/at_engine.h              # the whole public surface, three functions
  src/at_lines.c               # line assembly and classification
  src/at_urc.c                 # the unsolicited table and its dispatch
  src/at_backoff.c             # exponential with a cap and jitter
  src/modem_sim7020e.c         # the command set for this module, nothing else
  src/uart_dma.c               # from chapter 4, unchanged
  src/ringbuf.c                # from chapter 2, unchanged
  src/main.c
  test/test_engine.c           # host build, no hardware, chapter 12 runs it
  test/faults.py               # the hundred injected faults
  test/transcripts/            # recorded module replies, good and damaged
  docs/timeouts.md             # every timeout with the manual page it came from
  README.md
\end{plaincode}

The timeouts document is not decoration. It is the file a reviewer opens first,
because every number in it should be traceable to a page of the module's own
manual, and the ones that are not should be marked as measured on the bench with
the instrument that measured them.

\subsection*{Steps}

\begin{steps}
\item \textbf{Talk to the module from the host before writing any firmware.}
Power it from its own supply, connect the serial cable with the red 5 V lead
disconnected, and confirm it answers. This settles the baud rate, the line
ending the module expects and whether echo is on, all of which otherwise become
firmware bugs.
\begin{shellcode}
python -m serial.tools.miniterm /dev/ttyUSB0 115200 --eol CRLF
# then, one at a time:  AT   ATE0   AT+CMEE=2   AT+CPIN?   AT+CSQ
\end{shellcode}
Turning echo off with \texttt{ATE0} and turning verbose errors on with
\texttt{AT+CMEE=2} are the first two commands the engine will send, for the same
reasons you want them here: echo doubles every line the parser has to consider,
and a numeric error code with no text is a support request waiting to happen.

\item \textbf{Build the byte path and never parse inside it.} The receive side is
chapter 4's circular DMA with idle-line detection, feeding chapter 2's ring
buffer. The interrupt does one thing: it records how far the transfer counter
has moved and sets a flag. Line assembly happens in the engine, at whatever
priority the main loop runs at.

\item \textbf{Define the command record before writing the engine.} The record is
the contract, and most design questions answer themselves once it exists.
\begin{ccode}
typedef enum { AT_OK, AT_ERROR, AT_TIMEOUT, AT_ABANDONED } at_result_t;

typedef void (*at_done_fn)(at_result_t r, const char *last_line, void *ctx);

typedef struct {
    const char  *text;         /* without the line ending; the engine adds it */
    uint32_t     timeout_ms;   /* from the module's manual, never invented    */
    uint8_t      max_retries;  /* 0 means send once and report what happened  */
    uint8_t      flags;        /* AT_F_NO_ECHO_CHECK, AT_F_BINARY_TAIL, ...   */
    at_done_fn   done;         /* called once, in every outcome, including    */
    void        *ctx;          /* timeout and abandonment                     */
} at_cmd_t;
\end{ccode}
The completion callback is called exactly once for every command submitted, in
every outcome. That single rule removes a whole family of leaks in which the
caller waits for a callback that a failure path never reaches.

\item \textbf{Write the engine as a function that returns immediately.} Three
functions are the entire public surface: submit, poll and cancel.
\begin{ccode}
/* Called from the main loop as often as convenient. Never waits, never spins,
 * and performs at most one state transition per call so the worst-case time
 * through it is bounded and measurable.                                     */
void at_poll(at_engine_t *e, uint32_t now_ms)
{
    const char *line;
    while ((line = at_next_line(e)) != NULL) {     /* bounded: ring is finite */
        if (at_urc_dispatch(e, line))              /* unsolicited, any state  */
            continue;
        at_feed_reply(e, line);                    /* belongs to the command  */
    }
    if (e->state == AT_ST_WAITING && (int32_t)(now_ms - e->deadline_ms) >= 0)
        at_on_timeout(e, now_ms);
    if (e->state == AT_ST_IDLE)
        at_start_next(e, now_ms);
}
\end{ccode}
The deadline comparison is written as a signed difference on purpose, so that it
stays correct when the millisecond counter wraps. A comparison written the
obvious way is correct for forty-nine days and then is not.

\item \textbf{Classify every line before the state machine sees it.} Three
classes and nothing else: a final result code ends the command in flight, an
unsolicited code is dispatched from the table, and anything else is an
intermediate line belonging to the command in flight.
\begin{ccode}
/* Final result codes, per the standard command set. The module's own manual
 * adds module-specific ones and they go in this table, not in the code.    */
static const char *const FINAL_OK[]  = { "OK", "CONNECT", NULL };
static const char *const FINAL_BAD[] = { "ERROR", "+CME ERROR:", "+CMS ERROR:",
                                         "NO CARRIER", "ABORTED", NULL };

at_class_t at_classify(const char *line)
{
    if (match_any(line, FINAL_OK))  return AT_CLASS_FINAL_OK;
    if (match_any(line, FINAL_BAD)) return AT_CLASS_FINAL_BAD;
    if (at_urc_lookup(line))        return AT_CLASS_URC;
    return AT_CLASS_INTERMEDIATE;
}
\end{ccode}
An empty line is not a class. Modules pad replies with blank lines and a
classifier that treats one as an intermediate result will attach it to whatever
command is running. Skip them in the line assembler, before classification.

\item \textbf{Give the unsolicited codes a table and a handler each.} Registration
status, signal quality and incoming payload notifications all arrive without
being asked for, and some of them arrive in the middle of a command. The table
is data, the handlers are code, and neither touches the state machine.
\begin{ccode}
static const at_urc_t URC_TABLE[] = {
    { "+CEREG:",  on_registration },   /* network registration changed      */
    { "+CSQ:",    on_signal       },   /* only unsolicited if enabled       */
    { "RDY",      on_module_ready },   /* the module restarted underneath us */
    { NULL,       NULL }
};
\end{ccode}
The module-restart notification deserves its own handler even if it does nothing
but set a flag, because everything the engine believes about the module's state
is wrong after one, and the flag is what tells the application to reconfigure
rather than to carry on.

\item \textbf{Take every timeout from the module's manual and write down where.}
Some commands answer in milliseconds and some take tens of seconds. Network
registration is the long one, and the published empirical study of this exact
module is the source for what to expect in practice, with the caveat that its
host is a single-board computer and not this board.
\begin{ccode}
/* docs/timeouts.md carries the manual page for every one of these.          */
static const at_cmd_t CMD_ECHO_OFF = { "ATE0",     300,  1, 0, on_done, NULL };
static const at_cmd_t CMD_VERBOSE  = { "AT+CMEE=2", 300, 1, 0, on_done, NULL };
static const at_cmd_t CMD_SIM      = { "AT+CPIN?",  5000, 0, 0, on_sim,  NULL };
static const at_cmd_t CMD_SIGNAL   = { "AT+CSQ",    2000, 2, 0, on_csq,  NULL };
static const at_cmd_t CMD_REGSTAT  = { "AT+CEREG?", 5000, 2, 0, on_reg,  NULL };
\end{ccode}
The data-plane commands are module-specific and are not reproduced here from
memory. They come from the module's own AT command manual, which is the
normative reference for this chapter, and they go in one file so that porting to
the SIM7070G touches one file and not the engine.

\item \textbf{Back off, with a cap and with jitter.} Retrying immediately after a
failure that was caused by the network being busy is the one strategy guaranteed
to make things worse. Double the interval, cap it, and scatter it.
\begin{ccode}
uint32_t at_backoff_ms(uint8_t attempt, uint32_t base_ms, uint32_t cap_ms)
{
    uint32_t d = base_ms;
    for (uint8_t i = 0; i < attempt && d < cap_ms; i++) d <<= 1;
    if (d > cap_ms) d = cap_ms;
    return d - (d >> 3) + (at_rand_small() % (d >> 2 ? d >> 2 : 1));
}
\end{ccode}
The jitter matters more than it looks. Without it a fleet of nodes that all lose
the network at the same time all return at the same time.

\item \textbf{Inject a hundred faults and require a defined state from every
one.} The injector sits on the host, between the module and the board, or
replaces the module entirely with a recorded transcript. Four classes, and they
are the four that actually occur.
\begin{pycode}
FAULTS = [
    ("no_reply",      lambda c: b""),                     # module says nothing
    ("truncated",     lambda c: c.reply[: len(c.reply) // 2]),
    ("urc_midway",    lambda c: c.reply[:3] + b"\r\n+CEREG: 0\r\n" + c.reply[3:]),
    ("late_reply",    lambda c: (c.delay(c.timeout_ms + 500), c.reply)[1]),
]
# 25 of each, across the command set, with the engine's state logged after each.
\end{pycode}
The acceptance rule is not that the engine succeeds. It is that every one of the
hundred ends in one of the named results, that the completion callback fires
exactly once in each, and that none of them leaves the engine in a state from
which the next command cannot start.

\item \textbf{Run it with no SIM card and call the stop a pass.} With no card
fitted, the card status query returns an error and the engine has nowhere to go.
The correct behaviour is to stop in a named state, report it once and not retry
forever, because no amount of backoff produces a SIM card.
\begin{shellcode}
python test/faults.py --port /dev/ttyUSB0 --count 100 --report faults.json
python test/faults.py --scenario no_sim --expect-state AT_ST_STOPPED_NO_SIM
\end{shellcode}
\end{steps}

\subsection*{Build, flash and debug}

\diagram{c11_timing}{Two exchanges on the serial link. The first carries an unsolicited message arriving in the middle of a command, which a blocking driver would read as the reply. The second shows a reply that arrives after its deadline and is discarded rather than being attached to the next command.}

\begin{shellcode}
cmake -B build-fw -G Ninja -DCMAKE_TOOLCHAIN_FILE=cmake/arm-none-eabi.cmake
cmake --build build-fw -j
cp build-fw/firmware.bin "$PROBE_DISK"/   # onto the probe disk
\end{shellcode}

Debugging this chapter is mostly reading the link rather than the code. Keep a
transcript of every byte in both directions, with a timestamp, and write it to
the console on a second port so that reading it does not change the timing of
the thing being read. A transcript with timestamps answers almost every question
this chapter raises, and a transcript without them answers none of the
interesting ones.

\begin{note}{When the module answers the host but not the board}
In order of likelihood: transmit and receive are not crossed, which is the most
common wiring fault on this bench and produces exactly this symptom; the grounds
are not common, so the signalling has no reference and reads as noise; the
module is at a different baud rate because a previous experiment changed it; the
line ending the module expects is not the one being sent; echo is on and the
parser is treating the echoed command as the reply; or the module is resetting
on every transmit because it is being fed from the board rather than from its
own supply. Check them in that order and check the supply first if the symptom
comes and goes.
\end{note}

To recover from a module that has been left in an unknown configuration, use the
power control line rather than the command set. A module that will not answer at
any rate is easier to restart than to interrogate, and the restart notification
it emits when it comes back is already in the unsolicited table.

\subsection*{Verification and acceptance criteria}

\begin{itemize}
\item One hundred injected faults, twenty-five of each of the four classes,
every one ending in a defined result, with the completion callback fired exactly
once per submitted command and none of them leaving the engine unable to start
the next command.
\item No run reaches a state from which no transition is possible other than the
two named stopped states, and those two are reported as outcomes rather than
counted as failures.
\item With no SIM card fitted the engine stops in the named no-card state,
reports once, and does not retry, which is recorded as a pass.
\item The worst-case time spent inside the poll function, measured with the
cycle counter, stays below the budget with the longest reply the module can
produce present in the ring.
\item The longest interval with interrupts disabled stays below the budget,
measured the same way, because an engine that never blocks the main loop but
disables interrupts for a millisecond has moved the problem rather than solved
it.
\item Every timeout in the firmware appears in the timeouts document with either
the manual page it came from or the bench measurement and instrument that
produced it.
\item A late reply arriving after its deadline is discarded and is not attached
to the following command, proven by a transcript from the fault injector rather
than by inspection of the code.
\item The engine compiles and its tests run on the host with no hardware
present, which is what chapter 12 needs from it.
\end{itemize}

\subsection*{Variants}

\begin{table}[H]\centering\small
\begin{tabular}{p{24mm}p{26mm}p{44mm}p{22mm}p{20mm}}
\toprule
Axis & Variant & What changes & Cost & Built in full in \\
\midrule
Intelligence and reach & Edge host over a framed link & The board speaks a framed protocol to a Raspberry Pi and the Pi owns the modem, which removes the command set from the firmware entirely and is the right answer whenever a host is present & A host must exist and stay powered & Here and chapter 9 \\
Operating system & The maintained library on a kernel & With an operating system underneath, the MIT cellular library named above becomes the correct choice and this engine becomes the thing that explains it & A kernel, and its memory & Chapter 20 \\
Peripheral substitution & The SIM7070G instead & The same engine, a different command file, and a module the maintained library does support, which makes it the control case for this chapter & Different network and different manual & Chapter 17 \\
Peripheral substitution & The SIM7600E-H instead & A module that does offer a point-to-point protocol, so a network stack on the board becomes possible and the command set stops being the only route & A network stack, and far more code & Chapter 17 \\
Execution model & DMA receive with idle line & The byte path this chapter stands on & Setup complexity & Chapter 4 \\
Synchronisation & A queue between application and engine & With a kernel, the submit call becomes a queue send and the poll loop becomes a task & Kernel overhead & Chapter 20 \\
Language & Host tooling in Python & The fault injector here & None & Chapter 3 \\
\bottomrule
\end{tabular}
\caption{Variants for chapter 11. The edge-host row is the one that most often wins in practice, and this chapter builds it in full alongside chapter 9 precisely so that the comparison is honest rather than rhetorical.}
\end{table}

\subsection*{Pitfalls}

\begin{itemize}
\item Parsing inside the receive interrupt. It puts unbounded work at high
priority, and every timing number the chapter produces afterwards is about the
parser rather than about the system.
\item Treating an unsolicited message as the reply to the command in flight.
This is the failure that a bench with a good signal never reproduces and that a
field deployment reproduces on the first day.
\item Attaching a late reply to the next command. Once this happens the engine is
one command out of step and stays that way, and every subsequent result is
wrong in a way that looks like a module fault.
\item Comparing deadlines with a plain unsigned comparison. It is correct until
the millisecond counter wraps, which is far enough away to be missed and close
enough to happen.
\item Inventing timeouts. The module's manual states a maximum response time per
command and that is the number. A guessed timeout that is too long turns a
failure into a stopped system.
\item Retrying without backoff and without jitter, which converts one node's bad
minute into a fleet's bad hour.
\item Feeding the module from the board. The 2 A peak is real, and the symptom is
a module that resets whenever it tries to register, which looks like a firmware
fault and is not one.
\item Leaving echo on and writing the parser around it. Turn it off in the
second command the engine sends and the classifier gets simpler.
\item Allocating command records on a heap. An engine whose failure mode is heap
exhaustion has replaced one unbounded behaviour with another.
\end{itemize}

\subsection*{Best practices applied}

\begin{itemize}
\item The prior art is named precisely, including the one library that would
have done the job and the exact reason it cannot be used here. A chapter that
reimplements a maintained library without saying why has failed.
\item Every allocation is static and every buffer has a stated size and a stated
behaviour when it is full.
\item The completion callback fires exactly once per command in every outcome,
which is the single rule that keeps the caller's state machine honest.
\item Every timeout is traceable to a document or to a measurement with its
instrument named.
\item Failure is tested by injection rather than by waiting for it, and the
acceptance criterion is a defined state rather than a success.
\item A stop is allowed to be a pass. Engineering that cannot distinguish "this
cannot work and here is why" from "this is broken" produces systems that retry
forever.
\end{itemize}

\subsection*{Stretch goals}

\begin{itemize}
\item Record real transcripts from both modules on the bench and replay them
into the host build of the engine, so the test suite runs with no hardware and
still exercises real module behaviour.
\item Add a watchdog on the module rather than on the processor: if no line of
any kind has arrived for longer than the longest expected quiet period, restart
the module through its power control line and report it.
\item Price the engine in energy using chapter 10's ledger. The send phase stops
being a stub, and the registration interval becomes the most expensive decision
in the node.
\item Port the command file to the SIM7070G and run the same fault suite against
the maintained library on a kernel, which is a comparison nobody has published
for these modules.
\item Add a second queue at a higher priority for commands that must overtake a
long-running one, and measure what that does to the worst case through the poll
function.
\end{itemize}

\subsection*{Roadmap and next steps}

Chapter 12 takes this engine's host build and makes it a unit test suite that
runs on every commit, and then takes the board itself and makes it a
hardware-in-the-loop job. The fault injector written here is the most valuable
thing this chapter hands forward, because a test that can produce a late reply
on demand is worth more than any amount of careful reading.

For the wider path, the community embedded engineering roadmap places command
set work and connectivity where this chapter sits, between the peripheral
drivers below it and the system design above it, and the accompanying essay on
closing the gap between reading and building makes the case for exactly this
kind of exercise. The free course from a different silicon vendor on its own
connectivity software is the best structured material on the layer above this
one, and its concepts transfer even though its hardware does not. Appendix H
lists the courses. There is no course and no textbook covering this module on a
bare-metal microcontroller, which is why the chapter states its method in enough
detail to be repeated.

\subsection*{Portfolio evidence}

\begin{itemize}
\item A public repository with the engine, the module command file kept
separate from it, the host test build, the fault injector and the recorded
transcripts.
\item The fault report: one hundred injected faults, the class of each, the
result each produced, and the engine state after each, as a table rather than a
narrative.
\item The timeouts document, every entry traceable to a manual page or to a
bench measurement with its instrument named.
\item A transcript showing an unsolicited message arriving in the middle of a
command and the command still completing correctly, which is the clearest one
page summary of what the chapter is for.
\end{itemize}

\subsection*{Sources}

Normative references:
\begin{itemize}
\item The SIM7020 series AT command manual, for the command set, the final
result codes, the unsolicited result codes and the maximum response time of
every command used here. This is the normative reference for the chapter.
\item The SIM7020 series hardware design document, for the supply requirement,
the peak current and the power control line timing.
\item The 3GPP command set specification for terminal equipment, for the
standard commands and result codes that are not module specific.
\item Reference manual RM0455, for the USART, its idle-line detection and the
transfer engine that serves it. Not RM0433, which documents its sibling.
\item The board user manual for MB1363, for which USART instance reaches the
Arduino header pins.
\end{itemize}

Reusable implementations:
\begin{itemize}
\item Majerle's cellular library, MIT, which supports the SIM7070G on this
bench, bundles a messaging client, and requires an operating system.\newline
\url{https://github.com/MaJerle/lwcell}
\item Majerle's host-side parser for the radio coprocessor's command set, MIT,
whose queue and unsolicited dispatch structure this chapter follows.\newline
\url{https://github.com/MaJerle/lwesp}
\item Majerle's transfer-driven USART examples, MIT, which are the byte path
chapter 4 builds on and this chapter stands on.\newline
\url{https://github.com/MaJerle/stm32-usart-uart-dma-rx-tx}
\item The peer-reviewed empirical study of this exact module, for realistic
signal and registration expectations, noting its host is a single-board
computer.\newline
\url{https://jtde.telsoc.org/index.php/jtde/article/view/955}
\item A transport-agnostic messaging client, MIT, for the layer above this one
once a byte stream exists.\newline
\url{https://github.com/FreeRTOS/coreMQTT}
\end{itemize}
